% Pipeline diagram for EBC-SST paper (Item 13)
% Requires: \usepackage{tikz}
%           \usetikzlibrary{positioning, arrows.meta, calc}
%
% Usage: copy the tikzpicture block into a figure environment in sn-article-revised.tex
%
% =====================================================================
% NOTATION GUIDE — edit the strings in \def commands below to change
% any label without touching the diagram layout.
% =====================================================================

% -- Top row nodes (continuous / infinite-dimensional) --
\def\nodeTopA{$\mathcal{H}_d \cap \mathcal{C}_d(\lambda)$}     % curves in RKHS
\def\nodeTopB{$\mathcal{H}^*_d \times S^d_{++}$}               % dual reps x scale
%\def\nodeTopC{$? \times S^d_{++}$}                              % unknown quotient analog
%\def\nodeTopD{$?$}                                              % unknown learned analog

% -- Bottom row nodes (discrete / finite-dimensional) --
\def\nodeBotA{$\mathbb{R}^{n \times d}_*$}                     % landmark matrices
\def\nodeBotB{$St(d,n) \times S^d_{++}$}                       % Stiefel x SPD
\def\nodeBotC{$Gr(d,n) \times S^d_{++}$}                       % Grassmannian x SPD
\def\nodeBotD{$\mathcal{G}_r \times \mathcal{P}$}              % learned submanifold

% -- Horizontal arrow labels --
\def\labelFIE{\textrm{FIE}}           % Fredholm integral equation
\def\labelPiOne{$SVD$}              % SRQD + polar decomposition
\def\labelPiTwo{$\pi$}              % quotient projection St -> Gr
\def\labelTPCA{\textrm{tPCA}}         % tangent PCA

% -- Vertical / diagonal arrow labels --
\def\labelWPD{\textrm{SRQD}}           % weighted polar decomposition

% =====================================================================

\newlength{\rowsep}\setlength{\rowsep}{1.5cm}   % vertical gap between rows
\newlength{\colsep}\setlength{\colsep}{3.8cm}   % horizontal gap between columns

\begin{tikzpicture}[
    % --- node style ---
    every node/.style={font=\normalsize},
    obj/.style={inner sep=4pt},            % space around node text
    % --- arrow styles ---
    arr/.style={-{Stealth[length=5pt]}, thick},
    darr/.style={-{Stealth[length=5pt]}, thick, dashed},
    % --- label styles ---
    lbl/.style={font=\small, fill=white, inner sep=1.5pt},
]

% =====================================================================
% NODES — arranged on a 2 x 4 grid
% =====================================================================

% Top row  (y = 0)
\node[obj] (T1) at (0,0)            {\nodeTopA};
\node[obj] (T2) at (\colsep,0)      {\nodeTopB};
%\node[obj] (T3) at (2\colsep,0)     {\nodeTopC};
%\node[obj] (T4) at (3\colsep,0)     {\nodeTopD};

% Bottom row  (y = -\rowsep)
\node[obj] (B1) at (0,-\rowsep)         {\nodeBotA};
\node[obj] (B2) at (\colsep,-\rowsep)   {\nodeBotB};
\node[obj] (B3) at (2\colsep,-\rowsep)  {\nodeBotC};
\node[obj] (B4) at (3\colsep,-\rowsep)  {\nodeBotD};

% =====================================================================
% HORIZONTAL ARROWS
% =====================================================================

% Top row
\draw[arr]  (T1) -- node[lbl, above] {\labelFIE}  (T2);
% ? arrows omitted per author request — uncomment and label as needed:
\draw[darr]  (T2) -- node[lbl, above] {}  (B3);
% \draw[arr]  (T3) -- node[lbl, above] {$?$}  (T4);

% Bottom row
\draw[arr]  (B1) -- node[lbl, below] {\labelPiOne}   (B2);
\draw[arr]  (B2) -- node[lbl, below] {\labelPiTwo}   (B3);
\draw[arr]  (B3) -- node[lbl, below] {\labelTPCA}    (B4);

% =====================================================================
% VERTICAL / DIAGONAL ARROWS  (continuous <-> discrete)
% =====================================================================

% Column 1: discretization (dashed — approximation)
\draw[darr] (T1) -- (B1);

% Column 2: weighted polar decomposition (solid)
\draw[darr]  (T2) -- node[lbl, right] {\labelWPD}  (B2);

\end{tikzpicture}
